%% ma: L12-B08-0151
%% lop: 12
%% chuong: 2
%% bai: 8
%% dang: 12.8.51
%% loai: TN4
%% muc: TH
%% dap_an: "C"
%% nguon: "CĐ2.4.pdf giáo viên gửi 27/10/2026 (M07, Lê Minh Khôi – Nguyễn Anh Khôi, Chương II Chủ đề 4 Các bài toán về tâm tỉ cự), Câu 1"
%% kien_thuc: "Bài 8 (biểu thức tọa độ của phép toán vectơ, độ dài vectơ), Bài 6 (điểm thỏa hệ thức vectơ, quy tắc ba điểm)"
%% trang_thai: chinh-thuc
%% ghi_chu: "Giáo viên đã duyệt 27/10/2026"
\begin{ex}
Trong không gian $Oxyz$, cho hai điểm $P(1;4;-3)$ và $Q(-5;2;5)$. Tọa độ của điểm $M$ sao cho $\left|\overrightarrow{MP}+\overrightarrow{MQ}\right|$ đạt giá trị nhỏ nhất là
\choice{$M(-2;0;0)$}{$M(2;0;0)$}{\True $M(-2;3;1)$}{$M(-2;3;0)$}
\loigiai{Gọi $I$ là trung điểm $PQ$: $I(-2;3;1)$. $\left|\overrightarrow{MP}+\overrightarrow{MQ}\right|=2MI\ge0$, nhỏ nhất khi $M\equiv I(-2;3;1)$.}
\end{ex}
